With stability established, the frozen question became allocation:
\emph{can a local mechanism prevent first-arrival starvation of a later
distinct configuration?} Each registered experiment tightened the
answer and moved the identified bottleneck one step downstream. All
five share the identical identity gate (flag-off byte-identity
\texttt{9647ea8a0ca4dbd2}), the identical 12-run matrix (bac/bca/il/d
\(\times\) three seeds), and the identical primary endpoint:
second-block protected mass \(\ge 0.5 \times\) first-block protected
mass.

\subsection{Residual allocation gate (rule)}

The gate (protocol \cite{cllaprot} extension, record \cite{rule})
multiplies candidate accumulation by \(g_i = \max(0,1-R_i)\) with
headroom > 0, where \(R_i = I_{p,i}/(I_{p,i}+I_{w,i})\) is the
protected-share of the neuron's delivered input current. The
implementation behaves exactly as specified --- the gate sat at
\(g = 1.0\) through the entire second block (\(R\)-median 0.000),
headroom was reserved (0.249--0.294 vs.\ 0.155--0.161 in the no-rule
runs), and first-block permanence decelerated (170 events, 8.5/pres,
vs.\ 364, 18.2/pres) --- yet the primary endpoint failed 6/6 (ratios
0.008--0.192 vs.\ the 0.5 bar). The trajectory instrument located the
failure \emph{upstream of the rule}: the second-block working substrate
is destroyed during the first block (live unconsolidated C-channel
afferents \(197 \to 77\) under M2 normalization of the shrinking
working budget and M4 pruning), so at second-block onset
\(g=1\) with almost nothing to accumulate. Verdict: the gate is not the
bottleneck; substrate availability is \cite{rule}. Alternating
coexistence preserved (F1-il 3/3).

\subsection{Dormant-candidate reserve}

The reserve (record \cite{reserve}) lets ever-coactive candidates
survive dormancy at a floor instead of dying, protecting pre-associations
across inactive periods. The mechanism works measurably --- 90--255
reserved candidates per run survive --- but at second-block onset in the
blocked arms the reserved C-cohort (channels 8--15) is exactly
\emph{zero} (mean over the A block: reserved A-cohort 123.7, C-cohort
0.0); the block-2 permanence count is 38 vs.\ 35 without the reserve.
The record classifies the outcome as \emph{neither} preservation
failure \emph{nor} flux failure: a trigger keyed on ``ever-coactive''
is satisfiable only by the pattern that has already been active. The
reserve is vacuous for a genuinely first-seen pattern by the semantics
of its own trigger \cite{reserve}.

\subsection{First-exposure allocation}

The next amendment biases candidate redraw/eviction toward channels
that fired in the current window, so a currently active novel input can
bind candidate slots before churn eliminates its substrate (record
\cite{fe}). Primary endpoint still failed 6/6 (ratios 0.072--0.216),
but the causal path became exact: first C permanence at \(t = 47{,}400\)
--- 2.4 s after block onset, 10\(\times\) faster than baseline;
permanence events 41 vs.\ 35 (allocator-only) and 38 (reserve);
block-1 allocation \emph{increased} (248 vs.\ 170 events --- the
mechanism strengthens whatever is currently active). The binding supply
was capped by two interacting facts: only \(\approx 0.71\) visible
channels/neuron survived at onset (visibility cap), and the
redraw-bias branch was structurally inert (\texttt{connected()}
excludes exactly the visible channels it would have preferred). The
implementation audit (\cite{feimpl}) identified the latter as a spec
vacuity --- the bias could never fire --- and the next amendment
corrected the source.

\subsection{Source-selection correction}

The corrected source reads the module's per-window firing record
instead of afferent-derived channel sets, making all firing channels
observable regardless of synapse survival (record \cite{fec}). Supply
tripled (41 \(\to\) 61 C-cohort permanence events); first permanence at
2.5 s; primary endpoint still failed 6/6 (ratios 0.046--0.193).
The classification that survived into the paper's central question:
flux improved but remained asymmetric (3.05/pres C vs.\ 13.2/pres A),
and the measured residual bottleneck was \emph{post-permanence weight
growth + working-afferent churn} --- newly protected second-block
synapses consolidate at \(w\approx 0.02\) and grow little (mean 0.054
vs.\ the first block's 0.096) because the working C afferents that
would drive their LTP keep being destroyed (29 \(\to\) 11 across the
second block itself) by M2's working-mass squeeze. The record
registered the two candidate amendments it did not run (e1:
consolidation-time protection of a working-afferent fraction; e2:
protected-path LTP growth) \cite{fec}. Figure~\ref{fig:decomp} plots
the final state of this decomposition.

\subsection{What the sequence ruled out}

\begin{itemize}
\item Allocation gating: \textbf{ruled out as the limiter} --- the gate
did nothing wrong and changed nothing.
\item Dormant preservation: \textbf{ruled out} --- preservation works,
but has nothing preservable before first contact.
\item First-exposure visibility: \textbf{fixed then ruled out} --- the
source correction removed the vacuity; supply tripled.
\item Flux (candidate accumulation rate): \textbf{supported as
secondary} --- asymmetric (3.05 vs.\ 13.2/pres) but not the binding
constraint.
\item Working-substrate churn against the \emph{new} protected
structure: \textbf{the measured residual} --- leads directly to the
recruitment question (\S\ref{sec:recruitment}).
\end{itemize}